% Appendix C: Testing Documentation

\chapter{Testing Documentation}

\section{Testing Strategy}

The RM project implements a comprehensive testing strategy that encompasses unit testing, integration testing, end-to-end testing, and performance testing. The testing approach follows test-driven development (TDD) principles and ensures code quality throughout the development lifecycle.

\subsection{Testing Pyramid}

The project follows the testing pyramid approach with the following distribution:
\begin{itemize}
\item \textbf{Unit Tests (70\%):} Testing individual components and methods in isolation
\item \textbf{Integration Tests (20\%):} Testing component interactions and API endpoints
\item \textbf{End-to-End Tests (10\%):} Testing complete user workflows
\end{itemize}

\section{Backend Testing}

\subsection{Unit Testing with JUnit}

\subsubsection{Service Layer Testing}

\begin{lstlisting}
// UserServiceTest.java
@SpringBootTest
class UserServiceTest {

    @Autowired
    private UserService userService;

    @Autowired
    private UserRepository userRepository;

    @Test
    void createUser_shouldCreateUserSuccessfully() {
        // Given
        UserDTO userDTO = new UserDTO();
        userDTO.setLogin("testuser");
        userDTO.setEmail("test@example.com");
        userDTO.setPassword("password");

        // When
        UserDTO result = userService.createUser(userDTO);

        // Then
        assertThat(result.getId()).isNotNull();
        assertThat(result.getLogin()).isEqualTo("testuser");
        assertThat(result.getEmail()).isEqualTo("test@example.com");
    }

    @Test
    void createUser_withExistingLogin_shouldThrowException() {
        // Given
        UserDTO existingUser = createTestUser("existing");
        UserDTO duplicateUser = new UserDTO();
        duplicateUser.setLogin("existing");

        // When & Then
        assertThrows(BadRequestAlertException.class, () -> {
            userService.createUser(duplicateUser);
        });
    }
}
\end{lstlisting}

\subsubsection{Controller Testing with MockMvc}

\begin{lstlisting}
// AuthControllerTest.java
@SpringBootTest
@AutoConfigureMockMvc
class AuthControllerTest {

    @Autowired
    private MockMvc mockMvc;

    @Autowired
    private UserRepository userRepository;

    @Test
    void authenticate_withValidCredentials_shouldReturnToken() throws Exception {
        // Given
        createTestUser();

        // When
        mockMvc.perform(post("/api/authenticate")
                .contentType(MediaType.APPLICATION_JSON)
                .content("{\"username\":\"test\",\"password\":\"test\"}"))
                .andExpect(status().isOk())
                .andExpect(jsonPath("$.token").exists());
    }

    @Test
    void authenticate_withInvalidCredentials_shouldReturnUnauthorized() throws Exception {
        // When
        mockMvc.perform(post("/api/authenticate")
                .contentType(MediaType.APPLICATION_JSON)
                .content("{\"username\":\"invalid\",\"password\":\"invalid\"}"))
                .andExpect(status().isUnauthorized());
    }
}
\end{lstlisting}

\subsection{Integration Testing}

\subsubsection{Database Integration Tests}

\begin{lstlisting}
// RapportRequestRepositoryTest.java
@DataJpaTest
class RapportRequestRepositoryTest {

    @Autowired
    private RapportRequestRepository repository;

    @Autowired
    private TestEntityManager entityManager;

    @Test
    void findByTechnicianId_shouldReturnRequestsForTechnician() {
        // Given
        User technician = createTestUser();
        RapportRequest request = createTestRequest(technician);

        entityManager.persist(technician);
        entityManager.persist(request);
        entityManager.flush();

        // When
        List<RapportRequest> results = repository.findByTechnicianId(technician.getId());

        // Then
        assertThat(results).hasSize(1);
        assertThat(results.get(0).getId()).isEqualTo(request.getId());
    }
}
\end{lstlisting}

\section{Frontend Testing}

\subsection{Angular Component Testing}

\begin{lstlisting}
// login.component.spec.ts
describe('LoginComponent', () => {
  let component: LoginComponent;
  let fixture: ComponentFixture<LoginComponent>;
  let authServiceSpy: jasmine.SpyObj<AuthService>;

  beforeEach(async () => {
    const spy = jasmine.createSpyObj('AuthService', ['login', 'isAuthenticated']);

    await TestBed.configureTestingModule({
      declarations: [LoginComponent],
      providers: [
        { provide: AuthService, useValue: spy }
      ]
    }).compileComponents();

    fixture = TestBed.createComponent(LoginComponent);
    component = fixture.componentInstance;
    authServiceSpy = TestBed.inject(AuthService) as jasmine.SpyObj<AuthService>;
  });

  it('should create', () => {
    expect(component).toBeTruthy();
  });

  it('should call auth service login on form submit', () => {
    const loginResponse = { token: 'test-token' };
    authServiceSpy.login.and.returnValue(of(loginResponse));

    component.loginForm.setValue({
      username: 'testuser',
      password: 'testpass'
    });

    component.onSubmit();

    expect(authServiceSpy.login).toHaveBeenCalledWith({
      username: 'testuser',
      password: 'testpass'
    });
  });
});
\end{lstlisting}

\subsection{Angular Service Testing}

\begin{lstlisting}
// auth.service.spec.ts
describe('AuthService', () => {
  let service: AuthService;
  let httpMock: HttpTestingController;

  beforeEach(() => {
    TestBed.configureTestingModule({
      imports: [HttpClientTestingModule],
      providers: [AuthService]
    });

    service = TestBed.inject(AuthService);
    httpMock = TestBed.inject(HttpTestingController);
  });

  afterEach(() => {
    httpMock.verify();
  });

  it('should login and store token', () => {
    const mockResponse = { token: 'mock-jwt-token' };

    service.login({ username: 'test', password: 'pass' }).subscribe(response => {
      expect(response.token).toBe('mock-jwt-token');
      expect(localStorage.getItem('authToken')).toBe('mock-jwt-token');
    });

    const req = httpMock.expectOne('/api/authenticate');
    expect(req.request.method).toBe('POST');
    req.flush(mockResponse);
  });
});
\end{lstlisting}

\section{Mobile Testing}

\subsection{Flutter Widget Testing}

\begin{lstlisting}
// login_screen_test.dart
void main() {
  group('LoginScreen', () {
    late MockLoginBloc mockLoginBloc;

    setUp(() {
      mockLoginBloc = MockLoginBloc();
    });

    testWidgets('renders login form', (WidgetTester tester) async {
      when(() => mockLoginBloc.state).thenReturn(LoginInitial());

      await tester.pumpWidget(
        BlocProvider<LoginBloc>.value(
          value: mockLoginBloc,
          child: MaterialApp(home: LoginScreen()),
        ),
      );

      expect(find.byType(TextFormField), findsNWidgets(2)); // email and password
      expect(find.byType(ElevatedButton), findsOneWidget);
      expect(find.text('Login'), findsOneWidget);
    });

    testWidgets('shows loading indicator when logging in', (WidgetTester tester) async {
      when(() => mockLoginBloc.state).thenReturn(LoginLoading());

      await tester.pumpWidget(
        BlocProvider<LoginBloc>.value(
          value: mockLoginBloc,
          child: MaterialApp(home: LoginScreen()),
        ),
      );

      expect(find.byType(CircularProgressIndicator), findsOneWidget);
    });

    testWidgets('navigates to home on successful login', (WidgetTester tester) async {
      when(() => mockLoginBloc.state).thenReturn(LoginSuccess(user: mockUser));

      await tester.pumpWidget(
        BlocProvider<LoginBloc>.value(
          value: mockLoginBloc,
          child: MaterialApp(home: LoginScreen()),
        ),
      );

      // Verify navigation occurred
      // (This would require a mock navigator or integration test)
    });
  });
}
\end{lstlisting}

\subsection{Flutter BLoC Testing}

\begin{lstlisting}
// login_bloc_test.dart
void main() {
  group('LoginBloc', () {
    late MockLoginRepository mockRepository;
    late LoginBloc loginBloc;

    setUp(() {
      mockRepository = MockLoginRepository();
      loginBloc = LoginBloc(mockRepository);
    });

    tearDown(() {
      loginBloc.close();
    });

    test('initial state is LoginInitial', () {
      expect(loginBloc.state, equals(LoginInitial()));
    });

    blocTest<LoginBloc, LoginState>(
      'emits [LoginLoading, LoginSuccess] when login succeeds',
      build: () {
        when(() => mockRepository.login(any(), any()))
            .thenAnswer((_) async => mockLoginResponse);
        return loginBloc;
      },
      act: (bloc) => bloc.add(LoginRequested(email: 'test@test.com', password: 'password')),
      expect: () => [
        LoginLoading(),
        LoginSuccess(user: mockUser),
      ],
      verify: (_) {
        verify(() => mockRepository.login('test@test.com', 'password')).called(1);
      },
    );

    blocTest<LoginBloc, LoginState>(
      'emits [LoginLoading, LoginFailure] when login fails',
      build: () {
        when(() => mockRepository.login(any(), any()))
            .thenThrow(Exception('Login failed'));
        return loginBloc;
      },
      act: (bloc) => bloc.add(LoginRequested(email: 'test@test.com', password: 'password')),
      expect: () => [
        LoginLoading(),
        LoginFailure(error: 'Exception: Login failed'),
      ],
    );
  });
}
\end{lstlisting}

\section{End-to-End Testing}

\subsection{API Testing with Postman/Newman}

\subsubsection{Authentication Test Suite}

\begin{lstlisting}
// authentication-tests.postman_collection.json
{
  "info": {
    "name": "RM API Authentication Tests",
    "description": "End-to-end tests for authentication endpoints"
  },
  "item": [
    {
      "name": "User Registration",
      "request": {
        "method": "POST",
        "header": [
          {
            "key": "Content-Type",
            "value": "application/json"
          }
        ],
        "body": {
          "mode": "raw",
          "raw": "{\"login\":\"testuser\",\"email\":\"test@example.com\",\"password\":\"password\"}"
        },
        "url": {
          "raw": "{{baseUrl}}/api/register",
          "host": ["{{baseUrl}}"],
          "path": ["api", "register"]
        }
      },
      "event": [
        {
          "listen": "test",
          "script": {
            "exec": [
              "pm.test(\"Status code is 201\", function () {",
              "    pm.response.to.have.status(201);",
              "});",
              "",
              "pm.test(\"Response has user data\", function () {",
              "    var jsonData = pm.response.json();",
              "    pm.expect(jsonData).to.have.property('id');",
              "    pm.expect(jsonData.login).to.eql('testuser');",
              "});"
            ]
          }
        }
      ]
    }
  ]
}
\end{lstlisting}

\section{Performance Testing}

\subsection{JMeter Test Plan}

\subsubsection{Authentication Load Test}

\begin{lstlisting}
<!-- authentication-load-test.jmx -->
<?xml version="1.0" encoding="UTF-8"?>
<jmeterTestPlan version="1.2" properties="5.0" jmeter="5.4.1">
  <hashTree>
    <TestPlan guiclass="TestPlanGui" testclass="TestPlan" testname="RM Authentication Load Test" enabled="true">
      <elementProp name="TestPlan.user_defined_variables" elementType="Arguments" guiclass="ArgumentsPanel" testclass="Arguments" testname="User Defined Variables" enabled="true">
        <collectionProp name="Arguments.arguments">
          <elementProp name="baseUrl" elementType="Argument">
            <stringProp name="Argument.name">baseUrl</stringProp>
            <stringProp name="Argument.value">http://localhost:8080</stringProp>
          </elementProp>
        </collectionProp>
      </elementProp>
    </TestPlan>
    <hashTree>
      <ThreadGroup guiclass="ThreadGroupGui" testclass="ThreadGroup" testname="Authentication Load Test" enabled="true">
        <intProp name="ThreadGroup.num_threads">50</intProp>
        <intProp name="ThreadGroup.ramp_time">10</intProp>
        <longProp name="ThreadGroup.duration">300</longProp>
      </ThreadGroup>
      <hashTree>
        <HTTPSamplerProxy guiclass="HttpTestSampleGui" testclass="HTTPSamplerProxy" testname="Login Request" enabled="true">
          <stringProp name="HTTPSampler.domain">${baseUrl}</stringProp>
          <stringProp name="HTTPSampler.path">/api/authenticate</stringProp>
          <stringProp name="HTTPSampler.method">POST</stringProp>
          <boolProp name="HTTPSampler.follow_redirects">true</boolProp>
          <stringProp name="HTTPSampler.protocol">http</stringProp>
          <boolProp name="HTTPSampler.use_keepalive">true</boolProp>
          <stringProp name="HTTPSampler.postBodyRaw">true</stringProp>
          <elementProp name="HTTPsampler.Arguments" elementType="Arguments">
            <collectionProp name="Arguments.arguments">
              <elementProp name="" elementType="HTTPArgument">
                <boolProp name="HTTPArgument.always_encode">false</boolProp>
                <stringProp name="Argument.value">{"username":"testuser","password":"password"}</stringProp>
                <stringProp name="Argument.metadata">=</stringProp>
              </elementProp>
            </collectionProp>
          </elementProp>
        </HTTPSamplerProxy>
      </hashTree>
    </hashTree>
  </hashTree>
</jmeterTestPlan>
\end{lstlisting}

\section{Test Results and Coverage}

\subsection{Code Coverage Report}

\begin{table}[H]
\centering
\caption{Test Coverage Summary}
\label{tab:test_coverage}
\begin{tabular}{@{}lrrr@{}}
\toprule
Component & Lines (\%) & Branches (\%) & Functions (\%) \\
\midrule
Backend Services & 85.2 & 78.4 & 92.1 \\
Frontend Components & 82.6 & 75.8 & 88.9 \\
Mobile Widgets & 79.3 & 72.1 & 85.7 \\
Overall & 82.4 & 75.4 & 88.9 \\
\bottomrule
\end{tabular}
\end{table}

\subsection{Performance Test Results}

\begin{table}[H]
\centering
\caption{API Performance Benchmarks}
\label{tab:performance_results}
\begin{tabular}{@{}lrrrr@{}}
\toprule
Endpoint & Method & Avg Response (ms) & 95th Percentile (ms) & Error Rate (\%) \\
\midrule
/api/authenticate & POST & 245 & 412 & 0.1 \\
/api/rapport-requests & GET & 156 & 289 & 0.0 \\
/api/rapport-templates & POST & 334 & 567 & 0.2 \\
/api/files/upload & POST & 1245 & 2134 & 1.8 \\
\bottomrule
\end{tabular}
\end{table}

\section{Continuous Integration}

\subsection{GitHub Actions CI/CD Pipeline}

\begin{lstlisting}
# .github/workflows/ci.yml
name: CI/CD Pipeline

on:
  push:
    branches: [ main, develop ]
  pull_request:
    branches: [ main ]

jobs:
  test:
    runs-on: ubuntu-latest
    services:
      postgres:
        image: postgres:13
        env:
          POSTGRES_PASSWORD: postgres
        options: >-
          --health-cmd pg_isready
          --health-interval 10s
          --health-timeout 5s
          --health-retries 5

    steps:
    - uses: actions/checkout@v2
    
    - name: Set up JDK 11
      uses: actions/setup-java@v2
      with:
        java-version: '11'
        distribution: 'adopt'
    
    - name: Cache Maven packages
      uses: actions/cache@v2
      with:
        path: ~/.m2
        key: ${{ runner.os }}-m2-${{ hashFiles('**/pom.xml') }}
        restore-keys: ${{ runner.os }}-m2
    
    - name: Run tests
      run: mvn test
    
    - name: Generate coverage report
      run: mvn jacoco:report
    
    - name: Upload coverage to Codecov
      uses: codecov/codecov-action@v2
\end{lstlisting}